\section{Whole \(B_{16}\): geometry before the finite pair domain}
\label{app:arithwholeB16}

Now the ENTIRE incident forest is \(B_{16}\);
there is no spectator to delete or to charge.
Use actual unit roots
\(n_i=(e_i-e_{i+1})/\sqrt2\), \(i\le15\),
\(n_{16}=e_{16}\), and the original vector
\(z=E(p-x)\), \(x_1>\cdots>x_{16}>0\), \(\|x\|<1\).
The legal fresh bulk and B3/B2 terminal checks
give strong \(6,\ldots,14\) independently of five
and either endpoint.
We first prove all physical reductions needed
to enter the exact quadratic algorithm.

\subsection{The actual head and a tight short-edge length bound}
\label{app:arithB16edges}

For \(a=.643\) or \(.648\), the reference through
\((1,a),(6,f_6),\ldots,(14,f_{14})\) has positive
decreasing coordinates
\[
 a,\quad(f_6-a)/5\text{ five times},\quad
 \sqrt j-\sqrt{j-1}\ (7\le j\le14),\quad0,0 .
\]
Its first middle slope exceeds \(.274\), hence
the next square-root slope \(<.2\).
If actual \(x_1>a\), ONE primitive square gives
\[
 \|x\|^2>
 a^2+(f_6-a)^2/5+
             \sum_{i=6}^{13}(\sqrt{i+1}-\sqrt i)^2.
\]
The last sum is \(>.2117\).
A direct rational proof avoids a numerical receipt:
for \(A=4i+2\), put
\(L=(A^2-1)/(A(A^2-2))\).
The small root \(d=(\sqrt{i+1}-\sqrt i)^2\)
satisfies \(d(A-d)=1\), whereas
\[
                 L(A-L)-1=-{1\over A^2(A^2-2)^2}<0 .
\]
Thus \(d>L\).
For \(i=6,\ldots,13\), strict lower numerators
of \(L\) at denominator \(10^7\) are
\[
 385186,333704,294372,263340,238230,217494,200080,185248 .
\]
Their sum is \(2117654\).
With \(f_6>2.01815\), the two reference norms are
\(>1.0033565045\), \(>1.0070662045\).

The ACTUAL head's retained white diagram is \(B_{15}\).
Finite classification permits isolation or simple
three at its LONG endpoint two, with no symmetric
short-end alternative.
For the latter,
\(A_H^2=2g^2+2g\), \(C_H=2g+1\).
Own zero has zero time, own three has its affine
triangle, and parallel has the actual null pair.
Every remaining own has \(g\ge1/\sqrt2\), so
\(x_1>\sqrt{\sqrt2-1}>.643\), impossible by the
first reference.
The head is isolated:
\[
 c_H=A_Hp-\sqrt2g_He_1,\quad
 U=4g_H^2\in\mathcal O_K,\quad A_H^2=U/2-1,\quad U>2 .
\]
If \(U\ge12\), its ACTUAL edge has
\(C_H=\sqrt{U/2}\ge\sqrt6\), giving \(x_1>.648\),
impossible by the second reference.
The strict rational seal uses
\(\sqrt6>2.449\),
\(.648^2(3.449)<1.449\).
Thus \(2<U<12\).

Let \(C_S\) be the ACTUAL short-sixteen edge cosh.
Its actual finite endpoint and original comparison give
\(S_{16}\ge4\sqrt{(C_S-1)/(C_S+1)}\).
If \(C_S\ge16\), this is \(>3.7573\).
For strong five, interpolate through true knots
five through eight and this terminal pressure.
The last slope is \(<.172<\sqrt8-\sqrt7\).
Using
\[
 E_8>.69685+.0385186+.0333704=.768739,\qquad f_8<2.3971,
\]
its norm is
\[
                      >.768739+1.3602^2/8=1.000007005 .
\]
The upper seal for \(f_8\) follows from
\(\sqrt8<2.828428\), \(\sqrt2<1.414214\),
\(.431328(1.414214)<.61\); also \(f_8>2.39\)
gives the claimed last-slope bound.
For weak five use the actual pressures
\(x_1>.59\), \(S_6>2\) and \(S_{16}>3.7573\).
The decreasing slopes \(.59,.282,.17573\) give
one reference of norm
\[
                     .59^2+1.41^2/5+1.7573^2/10
                            =1.054530329>1 .
\]
Both alternatives contradict the SAME original
norm without adding a head gain.
Therefore \(C_S<16\).

\subsection{All seven actual short-wall profiles}
\label{app:arithB16shortwhite}

The retained ACTUAL \(A_{15}\) finite-white list is
isolation; helper three at \(j=1,2,14,15\);
or helper four at \(j=1,15\).
Two contacts give a cycle.
Other simple branches or non-three bonds cannot
give a finite rank-sixteen tree.
There are no foreign components.

Normalize old long roots to raw square two and
the genuine short root to square one:
\(\alpha_i=e_i-e_{i+1}\), \(s_{16}=e_{16}\).
All \(e_i\) are then \(K\)-valued.
A helper-three cap normalizes to square two,
helper product \(-1\), and short product \(-T/2\),
where \(T=2\sqrt2\gamma\).
The ACTUAL cycles give
\(T,T^2/2\in\mathcal O_K\), not \(T/2\) integral.
Its exact formulas are
\[
 \begin{gathered}
 c=Ap-\gamma P_{16}-P_j/\sqrt2,\quad
 2A^2=F_j(T)=4T^2+jT+j-2,\\
 C_S={8T+j\over\sqrt{32-j(16-j)}} .
 \end{gathered}
\]
The original old-plus-cap Gram is Lorentzian over
\(\mathbb Q(T)\), so actual \(K=\mathbb Q(T)\),
and every other \(F_j(T_\sigma)<0\).
A helper-four cap has raw square one and helper
product \(-1\), with
\[
       A^2=16\gamma^2+2j\gamma+j-1,\quad
       C_S=16\gamma+j,\quad K=\mathbb Q(\gamma).
\]
These field equalities follow from nonphysical
embeddings fixing an original indefinite principal,
not from arithmeticity of a face.

Own zero is retained before dispatch.
Helper-three at one/two has \(A^2=-1/2,0\).
At fifteen its actual null is
\(c+n_{14}+2n_{15}+\sqrt2n_{16}\);
at fourteen it is
\(c+n_{13}+2n_{14}+2n_{15}+\sqrt2n_{16}\).
Helper-four at one has zero time;
at fifteen use \(c+\sqrt2n_{15}+n_{16}\).
Isolation has \(A^2=-1\).
Every null has positive cap time and selected
products zero; ALL omitted original-wall signs
are nonpositive termwise.
Parallel own always has \(c+n_{16}\).

Near helper-four positive own has \(C_S\ge23\).
Far helper-four has \(\gamma<15/16\), so its
complete positive native alphabet is
three through eight.
The native-nine cubic at \(15/16\) is
\(-17/512\), so all larger finite labels, parallel
and dotted are beyond this cutoff.
Own three is rational and excluded by WHOLE-form
compact anisotropy/Meyer.
For owns four through eight use respectively
the nonphysical conjugates
\[
 -\sqrt2/2,\quad(1-\varphi)/2,\quad-\sqrt3/2,\quad
 \cos(3\pi/7),\quad-\cos(\pi/8).
\]
The resulting time squares are
\(8-\sqrt2\), \(9-5\varphi\), \(12-\sqrt3\),
\(16\gamma_\sigma^2+2\gamma_\sigma\), and
\(16\gamma_\sigma^2-2\gamma_\sigma\), all positive.
They contradict admissibility in the ACTUAL
generating field.

Far next-end-three has
\(T<15/4\), every other \(-1/2<T_\sigma<0\).
In degree \(d\ge2\), its nonzero integral norm
has magnitude \(<15/8<2\), hence one.
But \(T^2/2\) integral would make
\(\mathrm N(T)^2/2^d\) an integer, impossible.
Degree one is rational.
Near next-end-three has \(T<9/4\), leaving only
own three \(T=\sqrt2\) or own four \(T=2\).
Own five already has \(T^2=2\varphi^2>81/16\).
The rational case is impossible; the other
conjugate time numerator is \(20-14\sqrt2>0\).
Both next-end orientations and both helper-four
roles are exhausted.

For the remaining far-end-three put \(S=T\);
for near-end-three put \(S=T+2\).
Their time polynomials are \(4S^2+S-1\) and
\(4S^2-S-1\), with respective nonphysical intervals
\[
 { -1-\sqrt{17}\over8}<S_\sigma<{ -1+\sqrt{17}\over8},
 \qquad
 {1-\sqrt{17}\over8}<S_\sigma<{1+\sqrt{17}\over8}.
\]
Both have \(0<S<42/5\), \(|S_\sigma|<.641\)
and \(S^2/2\) integral.
For degree \(d\ge3\), the nonzero integer norm
is divisible by four but has magnitude
\(<(42/5)(.641)^2=3.4514004<4\).
Degree one is rational; actual degree is two.
Trace and norm are EVEN by the same half-square
argument already proved above.

The hand table uses the EXACT physical short-edge
cutoff \(C_S<16\), rather than only the coarser
\(S<42/5\) used to bound the degree:
\[
\begin{array}{c|c}
 \text{far}&\sqrt2\le S<(16\sqrt{17}-1)/8<65/8\\
 \text{near}&2+\sqrt2\le S<(16\sqrt{17}+1)/8<67/8 .
\end{array}
\]
Indeed the actual edge formulas are
\(C_S=(8S+1)/\sqrt{17}\) and
\(C_S=(8S-1)/\sqrt{17}\), respectively;
the lower bounds retain the native own gap.
For example the far integer \(4+3\sqrt2\)
passes the coarse bound but violates this
exact physical cutoff and is not a table row.

The far strips give traces \(2,4,6,8\),
\(-6<\mathrm N(S)<4\).
The near strips give traces \(4,6,8\),
\(-4<\mathrm N(S)<6\).
Solving ALL even trace/norm pairs and imposing
the two strict strips gives the COMPLETE eleven rows:
\[
\begin{array}{c|c|r}
 \text{role}&T&\mathrm N(F_j(T))\\ \hline
 \text{far}&2+\sqrt6&-53\\
 &3+\sqrt{13}&-25\\
 &3+\sqrt{11}&-151\\
 &3+\sqrt7&-19\\
 &4+\sqrt{14}&-117\\ \hline
 \text{near}&\sqrt2&-9\\
 &1+\sqrt{11}&-43\\
 &1+\sqrt7&-103\\
 &2+3\sqrt2&-137\\
 &2+\sqrt{14}&-229\\
 &2+2\sqrt3&-83 .
\end{array}
\]
For explicit coverage, the far norm columns are
\(-4,-2,2\).
Trace two gives only inadmissible conjugates or
nonreal roots; trace four retains \(2+\sqrt6\);
trace six retains all three \(3+\sqrt{13}\),
\(3+\sqrt{11}\), \(3+\sqrt7\);
trace eight retains only \(4+\sqrt{14}\).
The near \(S\) norm columns are \(-2,2,4\).
Trace four retains only \(2+\sqrt2\);
trace six retains \(3+\sqrt{11},3+\sqrt7\);
trace eight retains \(4+3\sqrt2,4+\sqrt{14},4+2\sqrt3\).
Subtract two to recover the printed near \(T\).
The omitted rows violate a strict conjugate/physical
bound or are repeated rational roots.
All displayed norms are negative and their physical
times positive, so the other time signs are exactly
the required negative ones.
These are necessary actual models, never examples.

\subsection{The original connected-basis squares remove every helper}
\label{app:arithB16helpersquares}

Normalize the actual head to the \(K\)-root
\(h=c_H/(\sqrt2g_H)\), with
\(h\cdot\alpha_1=-1\), \(h^2=2/U\).
Every nonphysical head Schur complement gives
\(0<U_\sigma<2\).
For any of the eleven helper-three rows, the
sixteen old raw roots and \(r=\sqrt2c\) are an
ORIGINAL \(K\)-basis.
The coefficient of \(r\) in \(h\), computed from
its positive time coefficients, therefore gives
\[
                      R_0={U-2\over U F_j(T)}\in K^2 .
\]
This derives a common time-line square in one
actual connected basis.  It does not put separate
unit roots or square roots in \(K\).

For quadratic \(K\), the full head strip gives
\[
                     0<\mathrm N(U)<24,\qquad
                     0<-\mathrm N(U-2)<20.
\]
For all five NEAR DOTTED rows and the far
\(\sqrt6,\sqrt{11}\) rows, the absolute time
norm is a prime at least forty-three.
Its valuation in \(\mathrm N(R_0)\) is odd
because neither of the two bounded integer norms
can contain it.
For far \(\sqrt7\), the FULL ring
\(\mathbb Z[\sqrt7]\) gives the complete head list
\[
 U=3+\sqrt7,\ 4+\sqrt7,\ 6+2\sqrt7,
\quad \mathrm N(U)=2,9,8,\quad
 \mathrm N(U-2)=-6,-3,-12 .
\]
Writing \(U=a+b\sqrt7\) shows \(b=1,2\);
the strict head strips give precisely these \(a\)'s.
None contains prime nineteen.
For far \(\sqrt{14}\), its full strip similarly
gives
\[
 U=4+\sqrt{14},\ 5+\sqrt{14},\quad
 \mathrm N(U)=2,11,\quad\mathrm N(U-2)=-10,-5 .
\]
The odd prime thirteen in time norm \(-117\)
therefore cannot cancel.
These are nine eliminated rows.

The near \(T=\sqrt2\) row needs its FULL square
test rather than a norm-only rejection.
Here actual \(K=\mathbb Q(\sqrt2)\), because its
connected old-plus-cap Gram is physically Lorentzian
over that field.
Put \(s=\sqrt2\), \(\epsilon=1+s\).
The ACTUAL cap has
\[
 A_F^2=(21+15s)/2=3\epsilon^3/2,\qquad
 c_F=A_Fp-P_{16}/2-P_{15}/s .
\]
The complete head strip in \(\mathbb Z[s]\) is
\[
         U=2+s,\ 3+s,\ 3+2s,\ 4+2s,\ 5+3s,\ 6+3s,\ 6+4s .
\]
Indeed \(U=a+bs\) has \(1\le b\le4\);
the allowed pairs are
\(b=1:a=2,3\);
\(b=2:a=3,4\);
\(b=3:a=5,6\);
\(b=4:a=6,7\), with the last \(a=7\) beyond twelve.
For \(R_0=(U-2)/(3\epsilon^3U)\), the norms are
\[
                      1/9,\ 1/63,\ 7/9,\ 1/18,\ 1/7,\ 1/81,\ 4/9 .
\]
The four nonsquares are impossible.
At \(U=2+s\), \(R_0=1/(3\epsilon^4)\);
at \(U=6+4s\), \(R_0=2/(3\epsilon^4)\).
Neither \(1/3\) nor \(2/3\) is a square in
\(\mathbb Q(s)\): a rational square in that field
requires the rational number itself or half of
it to be a rational square.
The last row \(U=6+3s\) REALLY gives a square,
with its positive physical root \(k=1/(3\epsilon)\).
For the ACTUAL mutual magnitude \(M=-c_H\cdot c_F\),
\[
 A_HA_F/g_H=\epsilon^2,\qquad
 M/g_H=(4+3s)/2,\qquad
                    4M^2=(174+123s)/2 .
\]
Its rational norm is \(9/2\), so it is not integral.
This contradicts the actual two-cycle.
The positive time sign is fixed by actual outward
roots, not an optional square-root sign.
Thus the near own-three model also is excluded.

The far \(T=3+\sqrt{13}\) model is last.
Put \(r=\sqrt{13}\), \(\eta=(3+r)/2\).
Its time polynomial is \(F_1(T)=5\eta^3\).
Use the FULL ring \(\mathbb Z[(1+r)/2]\).
The complete head strip gives the SIX rows
\[
\begin{array}{c|r|r|c}
 U&\mathrm N(U)&\mathrm N(U-2)&\mathrm N(R_0)\\ \hline
 (5+r)/2&3&-3&1/25\\
 (7+r)/2&9&-1&1/225\\
 4+r&3&-9&3/25\\
 5+r&12&-4&1/75\\
 (11+3r)/2&1&-17&17/25\\
 (13+3r)/2&13&-9&9/325 .
\end{array}
\]
Writing \(U=(a+br)/2\), same-parity integers,
the strip gives \(b=1,2,3\) with respectively
\(a=5,7;\ 8,10;\ 11,13\).
The last four norm entries are nonsquares.
For the first,
\(R_0=(101-28r)/15\).
A full field square would have rational coefficient
square \(52/15\) or \(49/15\), neither rational-square.
For the second, multiplying by the square \(\eta^2\)
gives
\[
                   \eta^2R_0=(7-r)/90 .
\]
Its rational coefficient square would be \(13/180\)
or \(1/180\), again impossible.
No norm-square is treated as sufficient.
All seven actual white profiles are now exhausted
except isolation.

\subsection{The actual isolated pair and the entire cubic branch}
\label{app:arithB16cubic}

The actual short cap is
\(c_S=A_Sp-g_SP_{16}\).
Own zero and parallel are impossible as above.
Put \(V=4g_S^2\).
Rescale the ACTUAL \(K\)-root to raw norm \(V\),
old-short product \(-V/2\):
physically \(F_S=\sqrt V\,c_S\), spatial
\((V/2)P_{16}\), time square \(V(4V-1)\).
Its original Lorentzian Gram over \(\mathbb Q(V)\)
gives ACTUAL \(K=\mathbb Q(V)\), and EVERY
nonphysical place has \(0<V_\sigma<1/4\).
The short edge bound gives \(V<64\).
Rational \(K\) is excluded by WHOLE-form Meyer.
Finite own would give \(V<4\) and
\(0<\mathrm N(V)<4(1/4)^{d-1}\le1\), impossible.
Thus \(V>4\), the short own is dotted.
Degree at least four gives
\(\mathrm N(V)<64(1/4)^3=1\).
Hence actual degree is two or three.

Normalize \(h=c_H/(\sqrt2g_H)\),
\(r_F=c_S/g_S\), so
\(h^2=2/U\), \(r_F^2=4/V\), spatial parts
\(e_1,P_{16}\).
The original old sixteen roots plus \(r_F\)
form a \(K\)-basis.
The positive coefficient of \(r_F\) in \(h\) gives
\[
                   k^2={(U-2)V\over4U(4V-1)}\in K^2,
                   \qquad k>0 .
\tag{B16K}\label{eq:arithB16square}
\]
Let \(H=Uh\), \(F=(V/2)r_F\).
The ACTUAL full long-path cycle and two-cycle give
\[
 \begin{gathered}
 G=-2H\cdot F=4Uk(4V-1)-UV\ge0,\qquad
 L={G^2\over2UV},\qquad G,L\in\mathcal O_K,\\
 (G+UV)^2=4UV(U-2)(4V-1),\\
 L+G=((15U-32)/2)V-2U+4 .
 \end{gathered}
\tag{B16C}\label{eq:arithB16cycles}
\]
The old-path factors are \(-1\), followed by
\(-V,-G/V\), whose complete product is \(G\).
Orthogonal \(G=L=0\) is retained.
The last identity, subtracting seven integral
copies of \(UV\), implies \(UV/2\) integral.
This is rank sixteen's own half-product, not
the odd-rank-seventeen census invariant.
Thus every actual pair has the COMPLETE domain
\[
 \begin{gathered}
 K=\mathbb Q(U)=\mathbb Q(V),\quad d=2\text{ or }3,\quad
 U,V,UV/2\in\mathcal O_K,\\
 2<U<12,\quad4<V<64,\qquad
 0<U_\sigma<2,\quad0<V_\sigma<1/4 .
 \end{gathered}
\]

Suppose \(d=3\).
The positive norm \(\mathrm N(V)\) is one, two or three.
For odd norm, \(V\) is a unit at every dyadic prime,
so \(UV/2\) integral makes \(U/2\) integral.
For norm two, \((V)\) is a SINGLE prime
\(\mathfrak p\) of norm two, exponent one.
Use integer prime valuations and let
\(e=\operatorname{ord}_{\mathfrak p}(2)\).
If \(U/2\) were nonintegral, its only deficient
prime could be \(\mathfrak p\), with
\(\operatorname{ord}_{\mathfrak p}(U)=e-1\).
Then
\(\operatorname{ord}_{\mathfrak p}(U-2)=e-1\),
\(\operatorname{ord}_{\mathfrak p}(4V-1)=0\).
The nonzero field square~\eqref{eq:arithB16square}
would have valuation
\[
                         (e-1)+1-2e-(e-1)=1-2e ,
\]
ODD, impossible for a square.
Thus \(W=U/2\) is integral in every case, with
physical \(1<W<6\) and both other conjugates
strictly between zero and one.

We prove the small cubic-integer fact needed here.
Any totally real cubic integer with physical root
\(w>1\) and other roots \(a,b\in(0,1)\) has \(w>5\).
Assume \(w\le5\).
Put integer trace \(S=w+a+b\le6\),
positive integer norms \(A=wab\), \(B=(w-1)(1-a)(1-b)\).
Its monic polynomial is
\[
              P(X)=X^3-SX^2+(A+S-B-1)X-A .
\]
It is positive between the two distinct small roots.
For \(0<X<1\), it increases with \(S\) and
decreases with \(A,B\).
If \(A\ge2\), then
\[
 P(X)\le X^3-6X^2+6X-2
                 <-5(X-3/5)^2-1/5<0 .
\]
If \(B\ge2\), similarly
\(P(X)<-5(X-2/5)^2-1/5<0\).
Thus \(A=B=1\).
For \(S\le5\),
\(P(X)\le X^3-5X^2+4X-1<-(2X-1)^2\le0\).
Only \(S=6\) remains, giving
\(P(X)=X^3-6X^2+5X-1\).
But \(P(5)=-1=(5-w)(5-a)(5-b)\) forces \(w>5\),
a contradiction.
Therefore \(U>10\).

The ACTUAL head edge now gives
\(x_1>(\sqrt5-1)/2>.618\).
The strict conjugate short bounds give
\(V>16\,\mathrm N(V)\ge16\), so its ACTUAL
short edge gives \(S_{16}>4\sqrt7/3>3.5276\).
Together with unchanged \(S_6>2.01815\), use
ONE positive reference
\[
 w_0=.33797P_1+.129085P_6+.150945P_{16}.
\]
Its coordinates are \(.618\), then five
\(.28003\)'s, then ten \(.150945\)'s.
The primitive square on the SAME original \(x\)
gives
\[
 \|x\|^2>
 .618^2+1.40015^2/5+1.50945^2/10
             =1.00185193475>1 .
\]
This excludes EVERY actual cubic field, with
all dyadic ramification allowed and no polynomial
field census.
Actual \(K\) is consequently quadratic.

\subsection{Exact inclusion in the distinct whole-\(B_{16}\) census}
\label{app:arithB16domain}

Every hypothetical whole \(B_{16}\) now satisfies
EXACTLY the actual-root strips and cycles of
Proposition~\ref{prop:pair-output}(iv), with the
original scales \(U=4g_H^2,V=4g_S^2\).
The square in~\eqref{eq:arithB16square} is a FULL
\(K\)-square, with its PHYSICAL positive root.
The exact \(G,L\) are
\eqref{eq:arithB16cycles}.
Acuteness gives \(G\ge0\).
For \(G=0\), \(L=L'=0\); otherwise strict
nonphysical two-root positivity gives \(0<L'<4\).
The own \(U\) and mutual \(L\) lie in the actual
quadratic finite native alphabet or are \(>4\).
Parallel four has its actual feasible null and
is excluded; zero mutual is retained.
These are necessary conditions of ACTUAL walls,
not conditions assigned to a formal projection.

For completeness, this algorithm has a finite
full-order domain justified by the physical strip.
For the maximal basis
\(\omega=(1+\sqrt D)/2\) or \(\sqrt D\),
\(U-U'=b\sqrt\Delta<12\), hence
fundamental discriminant \(\Delta<144\).
Trial \(2\le D\le143\), exact squarefreeness and
this discriminant test exhaust all actual fields.
The FULL safe integer boxes are
\[
 U=a+b\omega:\quad1\le a\le7,\ 1\le b\le5;\qquad
 V=a+b\omega:\quad1\le a\le32,\ 1\le b\le28 .
\]
For \(U\), the gap gives \(b\sqrt D<12\)
in the half basis and \(b\sqrt D<6\) otherwise;
\(0<U'<2\) bounds \(a\).
For \(V\), replace twelve/six by sixty-four/
thirty-two and use \(0<V'<1/4\).
These give the printed boxes, followed by EXACT
strict two-place strips.
At most ten head and twenty-eight short models
occur per field.
Visit EVERY ordered head/short pair.
Exact rational-pair arithmetic, both signed-norm
square-root branches, verified full square identities,
integrality, conjugate and native predicates are
the algorithm of Section~\ref{sec:algorithms}.
All loops are finite; there is no unbounded
coefficient search or numerical root cutoff.
The three disjoint discriminant intervals
\([0,48),[48,96),[96,144)\) cover the domain.
A separate complete combine recomputes field and
pair records; it is a check with the SAME
implementation, not an independent algorithm.

At this precisely proved domain we use the recorded
finite-output proposition~\ref{prop:pair-output}(iv)
as an EXPLICIT computed premise.
Its 43 fields,100 head models,85 short models
and273 pairs leave one necessary pair,
\[
                  U=5+2\sqrt6,\qquad V=10+4\sqrt6 .
\]
It is not an empty census or a realized polytope.
For this pair the ACTUAL head/short edges have
\[
 C_H=(2+\sqrt6)/2>11/5,\qquad
 C_S=2(2+\sqrt6)>44/5.
\]
They give \(x_1>.61\), \(S_{16}>3.56\).
With unchanged \(S_6>2.01815\), ONE reference
\[
                  w=.32837P_1+.127445P_6+.154185P_{16}
\]
has positive coefficients and norm
\[
                  .61^2+1.40815^2/5+1.54185^2/10
                              =1.00640742675>1 .
\]
The original completed square is contradictory.
Every actual short profile, all ground fields and
the complete nonempty quadratic pair domain have
now been covered for whole \(B_{16}\).
The finite output remains a declared numerical
premise; this source transcription is not a
new runtime replay or geometric Lean proof.

